% Part: first-order-logic
% Chapter: natural-deduction
% Section: provability-quantifiers

% verification of properties of provability needed for maximally
% consistent sets in the completeness chapter.

\documentclass[../../../include/open-logic-section]{subfiles}

\begin{document}

\olfileid{fol}{ntd}{qpr}

\olsection{\usetoken{S}{derivability} आणि संख्यापक}

\begin{explain}
  संपूर्णता प्रमेयासाठी या विभागात स्थापित केलेली~$\Proves$ संबंधी तथ्ये
  नैसर्गिक निगमनाच्या नियमांनी निष्पन्न होतात, हेही आवश्यक आहे.
\end{explain}

\begin{thm}
\ollabel{thm:strong-generalization} जर $c$ हे $\Gamma$ किंवा $!A(x)$ मध्ये
न येणारे स्थिर असेल आणि $\Gamma \Proves !A(c)$, तर $\Gamma
\Proves \lforall[x][!A(x)]$.
\end{thm}

\begin{proof}
$\delta$ ही $\Gamma$ पासून $!A(c)$ ची !!a{derivation} असू द्या.
\Intro{\lforall} अनुमान जोडल्यावर आपल्याला $\lforall[x][!A(x)]$ ची
!!a{derivation} मिळते. $c$ हे $\Gamma$ किंवा $!A(x)$ मध्ये येत नसल्यामुळे
आयगेन चराची अट पूर्ण होते.
\end{proof}

\begin{prop}
\ollabel{prop:provability-quantifiers}
\begin{tagenumerate}{prvEx,prvAll}
\tagitem{prvEx}{$!A(t) \Proves \lexists[x][!A(x)]$.}{}

\tagitem{prvAll}{$\lforall[x][!A(x)] \Proves !A(t)$.}{}
\end{tagenumerate}
\end{prop}

\begin{proof}
\begin{tagenumerate}{prvEx,prvAll}
\tagitem{prvEx}{पुढील ही~$\lexists[x][!A(x)]$ ची $!A(t)$ पासूनची
  !!a{derivation} आहे:
\begin{prooftree}
\AxiomC{$!A(t)$}
\RightLabel{\Intro{\lexists}}
\UnaryInfC{$\lexists[x][!A(x)]$}
\end{prooftree}}{}

\tagitem{prvAll}{पुढील ही~$!A(t)$ ची $\lforall[x][!A(x)]$ पासूनची
  !!a{derivation} आहे:
\begin{prooftree}
\AxiomC{$\lforall[x][!A(x)]$}
\RightLabel{\Elim{\lforall}}
\UnaryInfC{$!A(t)$}
\end{prooftree}}{}
\end{tagenumerate}
\end{proof}

\end{document}
